%% AMS-LaTeX Created with the Wolfram Language : www.wolfram.com

\documentclass[12pt]{article}
\usepackage[letterpaper,margin=1in]{geometry}
\usepackage{amsmath, amssymb, graphics, setspace}
\usepackage{graphicx} \graphicspath{ {./figs/} }
\usepackage{tcolorbox}
\usepackage[framed,numbered,autolinebreaks,useliterate]{mcode}
\newcommand{\mathsym}[1]{{}}
\newcommand{\unicode}[1]{{}}

\usepackage[english]{babel}
\newtheorem{theorem}{Theorem}

\usepackage{caption}
\usepackage{subcaption}

\usepackage{tabulary}
\usepackage{multirow}
\usepackage{lscape}

\usepackage{hyperref}
\hypersetup{
    colorlinks=true,
    linkcolor=blue,
    filecolor=magenta,      
    urlcolor=cyan,
    pdftitle={Overleaf Example},
    pdfpagemode=FullScreen,
    }

\def\mathbi#1{\textbf{\em #1}}

\begin{document}
\doublespacing	

\title{CE 401/5501: Applied AI \\ \large Topic 11: Supervised ML for Regression}
\author{Dr. ZhiQiang Chen}
\date{}
\maketitle

\section{Regression or Data-fitting Problems}
Given a dataset of $N$ points, where each point consists of an $m$-dimensional column vector $\mathbf{x} = (x_1, x_2, \dots, x_m)^T$ and a scalar target value $y$, which relates to the vector variable $\mathbf{x}$. As in our previous classification topics, the data set can be expressed as
\[
\mathcal{D} = \{ (\mathbf{x}_i, y_i) | i = 1, 2, \cdots, N\}
\]
Compared with the previous dataset for supervised learning, the target variable $y$ is not categorical, and in general, it is Real valued. 

\subsection{Problem Definition: Linear Regression}

In the following, we formulate the problem in the classic linear regression context. For this task, we want to fit the data with a linear model using the following form:

\begin{equation}
f(\mathbf{x}) = a_1 x_1 + a_2 x_2 + \cdots + a_m x_m + a_0
\label{p03t08eq01}
\end{equation}
which is a hyperplane defined in a space of $\mathbb{R}^{m+1}$, where $a_i$'s define the linear coefficients and $a_0$ is an intercept. 

To estimate the model, more specifically to estimate $a_i$'s from the data, the basic approach is called the least-square estimation method. In the following, we formulate the LSE using two ways: (1) scalar formulation; and (2) vectorized formulation. The former provides understanding to the latter.

\subsection{Scalar Formulation}
The LSE starts with defining the error of the prediction. The general principle for regression modeling is that an objective function is defined, which is often based on the errors, in terms of the modeling parameters. In this case we define $\mathbf{x}_j$ being the measured $j^{th}$ vector, $y_j$ the $j^{th}$ measured output, and $\bar{y}_j$ is the model prediction from the function model in Eq.~\ref{p03t08eq01}, $\bar{y}_j = f(\mathbf{x}_j)$. 

First, an individual error at the $j^{th}$ point is defined:
\[
\begin{split}
e_j & = y_j - \bar{y}_j \\
    & = y_j - f(\mathbf{x}_j)
\end{split}
\]

We can can tell this error can be either negative and positive. Fig.~\ref{part03topic08fig01} illustrates this assuming $\mathbf{x}$ is a 1-D variable.

\begin{figure}
\centering
\includegraphics[scale=0.5]{figs/part03topic08fig01.png}
\caption{Linear regression illustration.}
\label{part03topic08fig01}
\end{figure}


The question is to how to combine all errors from all data points? Apparently, if we keep the error expression above, errors can be cancelled out as they are either positive or negative. There are two ways to accumulate errors.
\begin{enumerate}
    \item [(1)] SAE: Sum up using the absolute of the errors.
    \item [(2)] SSE: Sum up using the squares of the errors. 
\end{enumerate}
Corresponding to the two schemes we have:
\begin{equation}
E1(\mathbf{a}) = \sum_{j = 1}^N |y_j - f(\mathbf{x}_j)| 
\label{p03t08eq02}
\end{equation}
\begin{equation}
E2(\mathbf{a})= \sum_{j = 1}^N |y_j - f(\mathbf{x}_j)|^2
\label{p03t08eq03}
\end{equation}
where $E1$ and $E2$ are functions of the model parameters $\mathbf{a} = (a_1, a_2, \cdots, a_m, a_0)$.

With the total errors defined above, they provide the objective function for estimating the regression model parameters $\mathbf{a}$. The common scheme is to minimize either the SAE or SSE above, which is to find $\mathbf{a}$ that minimizes the objective function. 
\begin{equation}
\hat{\mathbf{a}} = \arg_{\mathbf{a}} E(\mathbf{a})
\label{p03t08eq04}
\end{equation}

It is not straightforward, however, to tell which total error form (SSE or SAE) is better or which one has advantages over the other one. To illustrate this, we assume that only one data points are available ($N = 2$), and $\mathbf{x} = x$ being a 1-D vector variable. With these conditions, the total errors can be written as:
\[
E1(a_1, a_0) = |y_1 - a_1 x_1 - a_0|
\]
\[
\begin{split}
    E2(a_1, a_0) &= (y_1 - a_1 x_1 - a_0)^2\\
    & = y_1^2 - 2y_1a_1x_1 - 2y_1a_0 + a_1^2x_1^2 + 2a_1a_0*x_1 + a_0^2
\end{split}
\]
We see that the SSE can be expanded, which gives out a second-order or quadratic form in terms of the model parameters $a_1$ and $a_0$. Upon this, from our basic algebraic knowledge, we know that SSE can be minimized with a unique solution by taking the derivatives of $E2$ with respect to $a_1$ and $a_0$ respectivley, and set them equal to zero to solve for $a_1$ and $a_0$. Hence, $a_1$ and $a_0$ are solved directly (e.g., with using iterative procedures by starting with guess values to $a_1$ and $a_0$).

More formally, we state that the SSE as an objective function is convex (more specifically, it is quadratic). What about the SAE objective function $E1(a_1, a_0)$? We will state that it is still convex (as we can tell from plotting the function of $|a|$); however, there are derivatives (gradient) defined with respect to the vector of $(a_1, a_0)$ . A more specific programming algorithm needs to be used for solving this convex programming problem \footnote{Minimization of SAE is a linear programming problem, which can be solved using standard optimization techniques such as simplex method or interior point methods. A basic process of writing the linear constraints is involved in both methods.}. 

In the following, we will adopt the SSE approach to modeling the objective function. As a matter of fact, by minimizing an SSE objective, it is to obtain the least of squared errors, which is called a least-square (LS) problem. The process to estimate a model parameter is then called a least-square-estimation (LSE) process. 

By using the vectorized notations, we are meant to generalize the process above by considering $\mathbf{x}$ as a vector, and we will introduce a set of notations, and utilize vector calculus to estimate the modeling parameters directly due to the quadratic nature as illustrated above. 

\subsection{Vectorized Formulation: L-1 Norm and L-2 Norm}
The vectorized formulation starts with rewriting the linear model in Eq.~\ref{p03t08eq01}. In the following, we will take the following conversion: (1) Introduce $a_{m+1}$ and let $a_{m+1}$ as the intercept variable; and (2) introduce an augmented constant and treated as $x_{m+1}$; and (3) use the inner product operation by defining 
\[
\mathbf{a} = [a_1, a_2, \cdots, a_{m+1}]^T
\]
\[
\mathbf{x} = [x_1, x_2, \cdots, x_{m+1}]^T
\]

\begin{equation}
\begin{split}
f(\mathbf{x}) & = a_1 x_1 + a_2 x_2 + \cdots + a_m x_m + a_0 \\
              & = a_1 x_1 + a_2 x_2 + \cdots + a_m x_m + a_{m+1} \cdot 1 \\
              & = a_1 x_1 + a_2 x_2 + \cdots + a_m x_m + a_{m+1} x_{m+1} \\
              & = <\mathbf{x}, \mathbf{a}>
\end{split}
\label{p03t08eq05}
\end{equation}

With this nation, the previous total error definitions (as objective functions) in Eq.~\ref{p03t08eq02} and Eq.~\ref{p03t08eq03} become:

\begin{equation}
\begin{split}
E1(\mathbf{a}) & = \sum_{j = 1}^N |y_j - <\mathbf{x}_j, \mathbf{a}>| \\
               & = \sum_{j = 1}^N |y_j - \mathbf{x}_j^T \mathbf{a}|
\end{split}
\label{p03t08eq06}
\end{equation}

\begin{equation}
\begin{split}
E2(\mathbf{a}) & = \sum_{j = 1}^N |y_j - <\mathbf{x}_j, \mathbf{a}>|^2 \\
            &  = \sum_{j = 1}^N |y_j - \mathbf{x}_j^T \mathbf{a}|^2
\end{split}
\label{p03t08eq07}
\end{equation}
where $(y_j, \mathbf{x}_j)$ constructs the $j^{th}$ measured data pair. In addition, we introduce the matrix-vector product form of an inner product.

At this step, we introduce two notations that contain all $N$ measurements based the basic measurement $(y_j, \mathbf{x}_j)$:
\[
\mathbf{X} = \begin{bmatrix}
\mathbf{x}_1^T \\
\mathbf{x}_2^T \\
\vdots\\
\mathbf{x}_N^T    
\end{bmatrix}
\]
which is a $N \times (m+1)$ matrix \footnote{Note that in class, I used the transpose of $X$ above as $X$, which causes some inconsistency with the Everstine's text}. Note that this matrix is an augmented matrix, as the last row $\mathbf{x}_N^T$ is a one's row vector $[1, 1, \cdots, 1]$ as implied in Eq.~\ref{p03t08eq05}.


and 
\[
\mathbf{y} = \begin{bmatrix}
y_1 \\
\vdots\\
y_N 
\end{bmatrix}
\]
which is a $N \times 1$ vector.



By adopting the vectorized notations, we have two vectorized objective functions, which are termed (i) an L-1 Norm objective function, and (ii) a squared L-1 Norm objective function (Yes, a square operation is needed):
\begin{equation}
E1(\mathbf{a}) = \|\mathbf{y} - \mathbf{X} \mathbf{a} \|_1 \\
\label{p03t08eq08}
\end{equation}

\begin{equation}
E2(\mathbf{a}) = \|\mathbf{y} - \mathbf{X} \mathbf{a} \|_2^2
\label{p03t08eq09}
\end{equation}
We have obtained two compact vecterized form for the objective functions. In addition, the two forms correspond to the previous SAE and SSE formulation when using scalar variables.  



\subsection{LSE Solution}
Given the least-squares problem:

\begin{equation}
\min_{\mathbf{a}} E(\mathbf{a}) = \|\mathbf{y} - \mathbf{Xa} \|^2
\label{p03t08eq10}
\end{equation}

We compute the gradient of \(E(\mathbf{a})\) with respect to \(\mathbf{a}\) and then solve for \(\mathbf{a}\) by letting the gradient equal to a zero vector. 

First, rewrite the objective function:

\begin{equation}
\begin{split}
E(\mathbf{a}) &= \|\mathbf{y} - \mathbf{Xa} \|^2 \\
              &= (\mathbf{y} - \mathbf{Xa})^T(\mathbf{y} - \mathbf{Xa})
\end{split}
\label{p03t08eq11}
\end{equation}

Next, compute the gradient of \(E(\mathbf{a})\) with respect to \(\mathbf{a}\):

\begin{equation}
\begin{aligned}
\nabla E(\mathbf{a}) &= \frac{\partial E(\mathbf{a})}{\partial \mathbf{a}} \\
                     &= -2\mathbf{X}^T(\mathbf{y} - \mathbf{Xa})  \\
                     &= -2\left( \mathbf{X}^T \mathbf{y} -  \mathbf{X}^T \mathbf{Xa} \right)
\end{aligned}
\label{p03t08eq12}
\end{equation}

Now, set the gradient to zero and solve for \(\mathbf{a}\):

\begin{equation}
\begin{aligned}
\nabla E(\mathbf{a}) &= 0 \\
\mathbf{X}^T \mathbf{y} -  \mathbf{X}^T \mathbf{Xa}  & = 0 \\
\mathbf{X}^T \mathbf{Xa}  &= \mathbf{X}^T \mathbf{y}
\end{aligned}
\label{p03t08eq13}
\end{equation}

Finally, find the solution for \(\mathbf{a}\):

\begin{equation}
\begin{aligned}
\mathbf{a} &= (\mathbf{X}^T\mathbf{X})^{-1}\mathbf{X}^T\mathbf{y}
\end{aligned}
\label{p03t08eq14}
\end{equation}

At this moment, we should practice the inner-dimensionality rule for matrix-vector multiplication: $\mathbf{X}$ is a $N\times (m+1)$ matrix, $\mathbf{X}^T \mathbf{X}$ and its inverse are $(m+1) \times (m+1)$, $\mathbf{y}$ is $N\times 1$. As a result, $\mathbf{a}$ is expected as a $(m+1) \times 1$ column vector.






\section{Neural Network-Based Regression Using Multi-Layer Perceptron (MLP)}

\subsection{Introduction to MLP for Regression}
A Multi-Layer Perceptron (MLP) is a type of neural network that is well-suited for modeling nonlinear relationships due to its architecture, which includes multiple layers of neurons with nonlinear activation functions. This allows MLPs to learn complex patterns in the data, making them highly effective for regression tasks where the relationship between input variables and the output variable is not linear.

\subsection{Architecture of MLP}
An MLP for regression typically consists of:
\begin{itemize}
    \item An \textbf{input layer} that accepts feature vectors $\mathbf{x} = (x_1, x_2, \dots, x_m)^T$.
    \item One or more \textbf{hidden layers} where each neuron applies a non-linear activation function to the weighted sum of inputs from the previous layer.
    \item An \textbf{output layer} with a single neuron if the task is to predict a single continuous variable. The activation function in this layer might often be linear or none, depending on the nature of the output variable.
\end{itemize}

\subsection{Activation Functions}
Activation functions introduce non-linear properties to the network, enabling it to learn more complex decision boundaries. They are applied to the weighted sum of inputs:

\begin{equation}
a = \sigma(\mathbf{W}\mathbf{x} + \mathbf{b})
\end{equation}

where $\mathbf{W}$ and $\mathbf{b}$ are the weight matrix and bias vector of the layer, and $\sigma$ represents the activation function. Common choices of activation functions include:
\begin{itemize}
    \item \textbf{ReLU (Rectified Linear Unit):} $\sigma(x) = \max(0, x)$. It is popular due to its simplicity and effectiveness in alleviating the vanishing gradient problem.
    \item \textbf{Sigmoid:} $\sigma(x) = \frac{1}{1 + e^{-x}}$. It is useful for models where we need to predict probabilities as its output is always between 0 and 1.
    \item \textbf{Tanh (Hyperbolic Tangent):} $\sigma(x) = \tanh(x)$. This function outputs values between -1 and 1, making it more suitable for layers deep in the network where maintaining zero-centered data helps with learning efficiency.
\end{itemize}

For the output layer in regression tasks, a linear activation function (or no activation function) is often used to allow the model to output a range of values that are not bounded, which is suitable for regression predictions.

\subsection{Model Formulation}
The complete MLP model can be represented as a series of functional transformations:
\begin{equation}
f(\mathbf{x}) = \sigma^{(L)}(\mathbf{W}^{(L)} \sigma^{(L-1)}(\mathbf{W}^{(L-1)} \cdots \sigma^{(1)}(\mathbf{W}^{(1)}\mathbf{x} + \mathbf{b}^{(1)}) \cdots) + \mathbf{b}^{(L-1)}) + \mathbf{b}^{(L)})
\end{equation}

where $L$ is the number of layers, $\mathbf{W}^{(l)}$ and $\mathbf{b}^{(l)}$ are the weight matrix and bias vector for layer $l$, and $\sigma^{(l)}$ is the activation function used in layer $l$.

\subsection{Loss Function for Regression}
In regression tasks, the choice of the loss function is critical as it quantifies how well the predicted outputs match the actual target values. Commonly used loss functions in regression include:
\begin{itemize}
    \item \textbf{Mean Squared Error (MSE):} It is the average of the squares of the differences between the predicted and actual values.
    \begin{equation}
    \text{MSE} = \frac{1}{N} \sum_{i=1}^N (y_i - f(\mathbf{x}_i))^2
    \end{equation}
    \item \textbf{Mean Absolute Error (MAE):} It is the average of the absolute differences between predicted and actual values, providing a measure of error magnitude without considering direction.
    \begin{equation}
    \text{MAE} = \frac{1}{N} \sum_{i=1}^N |y_i - f(\mathbf{x}_i)|
    \end{equation}
    \item \textbf{Huber Loss:} A combination of MSE and MAE, which is quadratic for small errors and linear for large errors, defined by a parameter $\delta$.
    \begin{equation}
    L_\delta(y, f(\mathbf{x})) = \begin{cases} 
    \frac{1}{2}(y - f(\mathbf{x}))^2 & \text{for } |y - f(\mathbf{x})| \leq \delta \\
    \delta(|y - f(\mathbf{x})| - \frac{1}{2}\delta) & \text{otherwise}
    \end{cases}
    \end{equation}
\end{itemize}

These loss functions help in guiding the training process to minimize prediction errors effectively. The choice between these can depend on the specific characteristics of the data and the impact of outliers.


\subsection{Dangers of Overfitting}
Overfitting is a significant challenge in neural networks, particularly in complex models like MLPs with many layers or neurons. Overfitting occurs when a model learns the detail and noise in the training data to an extent that it negatively impacts the performance of the model on new data. Essentially, the model becomes too well-tailored to the training data, losing its ability to generalize.

\subsection{Strategies to Overcome Overfitting}
To prevent overfitting in MLPs, several strategies can be implemented:
\begin{itemize}
    \item \textbf{Regularization Techniques:} Adding regularization terms such as L1 or L2 penalties to the loss function can constrain the magnitude of the model weights and help prevent the model from becoming too complex.
    \item \textbf{Dropout:} Randomly dropping units in layers during the training phase can prevent complex co-adaptations on the training data. It effectively makes the model more robust and able to generalize better.
    \item \textbf{Early Stopping:} Monitoring the model's performance on a validation set during training and stopping the training process once the performance starts to degrade (i.e., when performance on the validation set worsens) can prevent overfitting.
    \item \textbf{Cross-Validation:} Using cross-validation techniques can help ensure that the model performs well across different subsets of the data and is not overfitting to a specific part of the training data.
    \item \textbf{Simplifying the Model:} Reducing the number of layers or the number of neurons in each layer can help make the model less complex and less prone to overfitting.
\end{itemize}

\subsection{Training MLP}
Training an MLP involves adjusting the network weights to minimize a loss function, typically using backpropagation and an optimization algorithm like stochastic gradient descent (SGD), Adam, or RMSprop. It is crucial to monitor not only the loss on the training data but also on a separate validation set to gauge the model's ability to generalize.







\section{Support Vector Regression (SVR)}

\subsection{Introduction to SVR}
Support Vector Regression (SVR) extends the concepts from Support Vector Machines (SVM) used in classification to regression problems. The goal of SVR is not only to fit the data but also to ensure that deviations from the actual values are within a certain threshold, maintaining a balance between model complexity and error minimization.

\subsection{Formulation of Linear SVR}
SVR constructs a linear model in a higher-dimensional feature space where the data points are mapped via a linear function. The formulation is similar to finding the "flatness" in SVM classification. The linear SVR function is modeled as:
\begin{equation}
f(\mathbf{x}) = \mathbf{w}^T \mathbf{x} + b
\end{equation}
where $\mathbf{w}$ is the weight vector and $b$ is the bias. The objective of SVR is to minimize the coefficient vector $\mathbf{w}$, which helps in reducing model complexity and overfitting.

\subsection{Objective Function with Epsilon-Insensitive Loss}
The epsilon-insensitive loss function used in SVR ignores errors that are within a certain distance $\epsilon$ of the true values, focusing on larger deviations which are more significant:
\begin{equation}
L_\epsilon(y, f(\mathbf{x})) = \max(0, |y - f(\mathbf{x})| - \epsilon)
\end{equation}
The complete objective function to be minimized under the constraints imposed by the data points becomes:
\begin{equation}
\min_{\mathbf{w}, b, \xi, \xi^*} \frac{1}{2} \|\mathbf{w}\|^2 + C \sum_{i=1}^N (\xi_i + \xi_i^*)
\end{equation}
subject to:
\begin{align*}
y_i - (\mathbf{w}^T \mathbf{x}_i + b) &\leq \epsilon + \xi_i, \\
(\mathbf{w}^T \mathbf{x}_i + b) - y_i &\leq \epsilon + \xi_i^*, \\
\xi_i, \xi_i^* &\geq 0, \forall i.
\end{align*}
where $\xi_i$ and $\xi_i^*$ are slack variables that measure the degree of misfit.

\subsection{Kernel Trick for Nonlinear SVR}
To handle non-linear relationships, SVR can be extended using the kernel trick, allowing the algorithm to perform linear regression in a transformed feature space:
\begin{equation}
k(\mathbf{x}_i, \mathbf{x}_j) = \langle \phi(\mathbf{x}_i), \phi(\mathbf{x}_j) \rangle
\end{equation}
where $\phi$ represents a transformation to a higher-dimensional space. Common kernels include the polynomial, radial basis function (RBF), and sigmoid kernels.

\subsection{Sparsity and Model Complexity}
In SVR, sparsity refers to the model relying on only a subset of the training data to form the decision function. Specifically, only support vectors, which are data points that lie outside the $\epsilon$ margin, contribute to the model. This leads to a model that is both simpler and often more robust to overfitting. The degree of sparsity is controlled by the parameter $C$, which balances the model's complexity against its performance on the training data.


\subsection{Optimization and Solution}
The dual form of the optimization problem is often solved using quadratic programming techniques that handle the constraints and optimize the margins. This dual formulation allows for the incorporation of kernels and solves the regression problem by maximizing the margin while minimizing the loss within the epsilon-insensitive zone.





\end{document}